% =============================================================================
% Architecture de la solution — TicketManager
% Exporter les .drawio en PNG dans docs/release-1/img/ :
%   - architecture-applicative-3-tiers.png
%   - architecture-physique.png
% =============================================================================

\section{Architecture de la solution}

Cette section décrit l'organisation globale de la solution \textit{TicketManager},
en distinguant l'architecture logicielle (organisation des couches applicatives)
et l'architecture physique (environnement de déploiement on-premises).

% ---------------------------------------------------------------------------
\subsection{Architecture logicielle}
% ---------------------------------------------------------------------------

\textit{TicketManager} est une application web full-stack basée sur une architecture
en trois tiers assurant la séparation entre la présentation, la logique métier
et la gestion des données. La figure~\ref{fig:archi-logicielle} présente cette
organisation sous une forme conceptuelle, indépendante des technologies
d'implémentation.

\begin{figure}[htbp]
    \centering
    \includegraphics[width=0.92\textwidth]{img/architecture-applicative-3-tiers.png}
    \caption{Architecture logicielle de TicketManager — Vue en trois couches}
    \label{fig:archi-logicielle}
\end{figure}

\begin{itemize}
    \item \textbf{Couche présentation} : fournit les interfaces utilisateur
    dédiées aux trois profils (Client, Technicien, Administrateur), la
    consultation des notifications et les tableaux de bord statistiques.
    Elle communique avec la couche métier via des API REST et des connexions
    WebSocket pour les alertes temps réel.

    \item \textbf{Couche métier} : centralise les traitements applicatifs :
    authentification et contrôle d'accès (RBAC), gestion du cycle de vie des
    tickets, administration des utilisateurs, diffusion des notifications et
    calcul des statistiques.

    \item \textbf{Couche données} : assure la persistance des informations
    dans une base de données documentaire et délègue le stockage des pièces
    jointes à un service externe spécialisé.
\end{itemize}

\paragraph{Technologies d'implémentation.}
La couche présentation est développée avec \textbf{Next.js} et \textbf{React},
la couche métier avec \textbf{NestJS}, et la couche données s'appuie sur
\textbf{MongoDB} (via Mongoose) complété par \textbf{Cloudinary} pour les
fichiers joints.

% ---------------------------------------------------------------------------
\subsection{Architecture physique}
% ---------------------------------------------------------------------------

L'architecture physique décrit l'environnement de déploiement de la solution
sur \textbf{Azure Stack Hub}, hébergée on-premises au sein de RFC Tunisie.
La figure~\ref{fig:architecture-physique} en donne une vue d'ensemble.

\begin{figure}[htbp]
    \centering
    \includegraphics[width=\textwidth]{img/architecture-physique.png}
    \caption{Vue d'ensemble de l'architecture physique de TicketManager}
    \label{fig:architecture-physique}
\end{figure}

L'infrastructure repose sur les composants suivants :

\begin{itemize}
    \item \textbf{Azure Stack Hub} : plateforme cloud hybride on-premises
    (groupe de ressources, réseau virtuel \texttt{10.1.0.0/16}).

    \item \textbf{Réseau virtuel segmenté} en trois sous-réseaux, chacun
    protégé par un groupe de sécurité réseau (NSG) :
    \begin{itemize}
        \item \texttt{mgmt} (\texttt{10.1.1.0/24}) — administration ;
        \item \texttt{db} (\texttt{10.1.2.0/24}) — base de données (privé) ;
        \item \texttt{aks} (\texttt{10.1.3.0/24}) — cluster Kubernetes.
    \end{itemize}

    \item \textbf{Cluster Kubernetes (AKS Engine)} : orchestre les conteneurs
    applicatifs (interface web, service API, base de données in-cluster) et les
    outils DevOps (analyse qualité, registre d'images, supervision, scan de
    sécurité). Le plan de contrôle (nœud master) et les nœuds worker hébergent
    les pods du namespace \texttt{ticketmanager}.

    \item \textbf{VM-gestion} (\texttt{10.1.1.4}) : machine d'administration
    avec adresse IP publique et privée ; héberge Ansible, Docker, AKS Engine
    et l'agent Azure DevOps pour l'automatisation et le déploiement du cluster.

    \item \textbf{VM-db} (\texttt{10.1.2.4}) : machine dédiée à la base de
    données, accessible uniquement en IP privée depuis les sous-réseaux
    \texttt{mgmt} et \texttt{aks}. Elle sert de support \textit{legacy},
    de sauvegarde ou de migration ; la base de données de production est
    hébergée in-cluster (StatefulSet avec volume persistant).

    \item \textbf{Chaîne CI/CD} : dépôt de code (GitHub / Azure Repos),
    pipeline Azure DevOps (build, tests, qualité, images, scan, déploiement)
    et provisionnement Terraform (\textit{Infrastructure as Code}).

    \item \textbf{Services externes} : stockage des pièces jointes et accès
    utilisateur via Internet (adresse IP publique exposant l'interface web).
\end{itemize}
